\begin{pastexam}[label={px:ncu-110-dsa-09}]{中央資工所 110 學年度 複選題第 9 題}
（複選；每題 5 分，全對才給分，答錯一個選項扣 1.5 分）

The pseudo code below aims to reverse a linked list.
\begin{minted}[fontsize=\small]{text}
typedef struct Node{
    int data;
    Node *next;
};
Node *NPtr;

int main()
{
    NPtr head = NULL; /* Start with an empty linked list */
    push(&head, ...); /* insert items into the linked list */
    reverse(&head); /* reverse the linked list */
}

void reverse(NPtr *Href)
{
    NPtr P = NULL;
    NPtr C = *Href;
    NPtr N = NULL;

    while (____(1)____) {
        N = C->next;
        C->next = P;
        P = C;
        C = N;
    }
    ____(2)____
}
\end{minted}
Which of the following statements are true?
\begin{enumerate}[label=(\Alph*)]
\item Blank (1) should be C != NULL.
\item Statement C->next = P; reverses node's pointer.
\item Blank (2) should be *Href = C;
\item Statement P = C; moves the pointer to the next.
\end{enumerate}
\end{pastexam}

\begin{pastsoln}
答案：(A)(B)(D)。

這是資結書虛擬碼~5.12 的三指標反轉：$P$ 是已經反轉好的那一段的頭，$C$ 是目前要處理的節點，
$N$ 先記住 $C$ 原本的下一個，免得改了 \texttt{C->next} 之後找不到後面。 程式碼本身是「虛擬碼」，
\texttt{Node *NPtr;} 應該是 \texttt{typedef Node *NPtr;} 的意思，不影響作答。

以串列 $1 \to 2 \to 3$ 追蹤，每一列是做完一圈迴圈之後：
\begin{center}
\small
\begin{tabular}{cccl}
\toprule
圈 & $P$ & $C$ & 已反轉的部分（從 $P$ 走） \\
\midrule
開始 & \NIL & $1$ & （空） \\
1 & $1$ & $2$ & $1 \to \NIL$ \\
2 & $2$ & $3$ & $2 \to 1 \to \NIL$ \\
3 & $3$ & \NIL & $3 \to 2 \to 1 \to \NIL$ \\
\bottomrule
\end{tabular}
\end{center}

\textbf{(A) 對}。 每一圈處理一個節點 $C$，處理到最後一個節點之後 $C$ 變成 \NIL，迴圈要停；
空串列時 $C$ 一開始就是 \NIL，一圈都不跑。 所以條件是 \texttt{C != NULL}。 寫成
\texttt{C->next != NULL} 會漏掉最後一個節點，空串列時還會存取空指標。

\textbf{(B) 對}。 \texttt{C->next = P;} 把 $C$ 的指標從「指向原本的下一個」改成「指向前一個」，
這一行就是反轉的動作。

\textbf{(C) 錯}。 迴圈結束時 $C = \NIL$，新的頭是原本的最後一個節點，存在 $P$ 裡（上表最後一列）。
空格 (2) 應該是 \texttt{*Href = P;}，寫 \texttt{*Href = C;} 會讓串列變成空的。

\textbf{(D) 對}。 每一圈最後兩行把兩個指標各往前推一格：\texttt{P = C;} 讓 $P$ 移到下一個節點
（剛反轉好的 $C$），\texttt{C = N;} 讓 $C$ 移到下一個節點。 選項說 \texttt{P = C;} 把指標往下一個移，
說的是 $P$，成立。
\end{pastsoln}
